\section{Arithmetic comparison maps}
\label{sec:arithmetic-comparison}

Proposition~\ref{prop:divisor-bar} uses the multiplication of effective
divisors and its augmentation module. To compare it with arithmetic cochains,
one must specify the coefficient category and the boundary modules there.  Let \((\mathcal C,\widehat\otimes,\one)\) be a stable, closed,
\(\mathbb C\)-linear symmetric monoidal category with a dg model.  Assume
that it admits geometric realizations and filtered colimits and that
tensoring preserves them separately in each variable. Here $E_1$ means
associative multiplication with its coherent higher homotopies. For an
\(E_1\)-algebra \(A\), a right \(A\)-module \(M\), and a left
\(A\)-module \(N\), the simplicial two-sided bar has
\[
 \operatorname{Bar}_A(M,N)_n
 =M\widehat\otimes A^{\widehat\otimes n}\widehat\otimes N.
\]
Its faces use the multiplication and module actions, while its degeneracies
use the unit.  Its realization computes \(M\widehat\otimes_A N\) under the
stated hypotheses
\cite[Construction~4.4.2.7 and Theorem~4.4.2.8]{LurieHA}.  An augmentation
\(\epsilon:A\to\one\) makes \(\one\) into both boundary modules and hence
defines the reduced bar \(B(A)=\operatorname{Bar}_A(\one,\one)\).

An associative chiral duality must compare this multiplication
and its bar with another associative object in a specified category of
factorization data. A Lie--commutative chiral duality has different input
and output operations and supplies no such comparison without an
additional functor and compatibility maps.

\begin{proposition}[One heat trace, inequivalent bars]
\label{prop:two-bars}
Fix \(p\) and let
\(V_p=\bigoplus_{r\geq0}\mathbb C e_r\), with
\(N_pe_r=r\log p\,e_r\).  Give $e_r$ weight $r$ and put every $e_r$ in homological degree zero.
The same pointed weight-graded vector space carries the
two augmented associative algebra structures
\[
 A_p^{\mathrm{pol}}=\mathbb C[u],\quad e_r=u^r,
 \qquad
 A_p^{\mathrm{sq}}=\mathbb C\oplus I,\quad
 I=\bigoplus_{r\geq1}\mathbb C e_r,\quad I^2=0,
\]
with unit \(e_0\) and augmentation \(e_0\mapsto1\),
\(e_r\mapsto0\) for \(r>0\).  In both cases \(N_p\) is a derivation, and
on the common Hilbert completion both algebras give the heat trace
\[
 \operatorname{Tr}(e^{-sN_p})=(1-p^{-s})^{-1},
 \qquad \operatorname{Re}s>0.
\]
Nevertheless, their reduced bars are not quasi-isomorphic.  For the
augmentation module \(\mathbb C\),
\[
 \operatorname{Tor}^{A_p^{\mathrm{pol}}}_n(\mathbb C,\mathbb C)
 =\begin{cases}
   \mathbb C,&n=0,\\
   \mathbb C\eta_p,&n=1,\\
   0,&n\geq2,
  \end{cases}
 \qquad
 \operatorname{Tor}^{A_p^{\mathrm{sq}}}_n(\mathbb C,\mathbb C)
 =I^{\otimes n}\quad(n\geq0).
\]
Thus the energy spectrum and its trace do not select an \(E_1\)-algebra or
a bar construction.
\end{proposition}

\begin{proof}
The two-term Koszul resolution of \(\mathbb C\) over \(\mathbb C[u]\)
gives the first calculation.  In the normalized bar of
\(\mathbb C\oplus I\), every multiplication face on the augmentation ideal
vanishes because \(I^2=0\), so the reduced bar differential is zero and its
degree-\(n\) term is \(I^{\otimes n}\).  The derivation identity is immediate
for the polynomial product and is still valid when two positive-weight
elements have zero product.  The heat trace depends only on the common
diagonal operator.
\end{proof}

\begin{remark}
The Toeplitz algebra does not remove this ambiguity.  It has no augmentation
sending \(S_p\) to zero, since \(S_p^*S_p=1\).  Every character sends
\(S_p\) to a scalar of modulus one, and no such character is invariant under
\(S_p\mapsto p^{it}S_p\).  A Toeplitz bar therefore requires explicitly
chosen boundary modules; it cannot silently inherit the polynomial
augmentation.
\end{remark}

\begin{openproblem}[Arithmetic-chiral realization]
\label{problem:bar}
Extend the divisor construction of Section~\ref{sec:divisor-bar} to
arithmetic local coefficients by constructing the following data and
compatibilities.  In particular, identify the maps from the divisor
algebras or their bars to the proposed arithmetic objects.

\begin{enumerate}
\item Fix one category
\((\mathcal C,\widehat\otimes,\one)\) as above, including the completion
functor or tensor topology.  For every place \(v\), give an augmented
\(E_1\)-algebra \((A_v,\eta_v,\epsilon_v)\), or give explicit right and
left boundary modules when no augmentation is intended.  At infinity,
also give the realization functor from the chosen chiral or factorization
algebra to \(\mathcal C\).

\item Construct each local bar through simplicial length three, with its
grading and signs.  For a finite set \(S\) of places containing infinity,
form the declared completed tensor product \(B_S\) of the local bars.

\item Use cohomological indexing for this twisting construction, so the
homological bar above is reindexed as \(B^{-n}=B_n\), with the same
differential.  Give a coaugmented conilpotent dg coalgebra
\[
 P_S\in\operatorname{CoAlg}_{\mathrm{conil}}(\mathcal C),
 \qquad \eta_{P_S}:\one\longrightarrow P_S,
\]
set
\(E_S=\underline{\operatorname{End}}_{\mathcal C}(B_S)\), and construct a
continuous degree-one map
\[
 \tau_S:P_S\longrightarrow E_S.
\]
Require the normalization \(\tau_S\circ\eta_{P_S}=0\).  Equip these
objects with complete separated descending filtrations for which the
differentials, coproduct, composition, and action on \(B_S\) are continuous.
Require the resulting twisting summand \(T_{\tau_S}\) either to be locally
nilpotent before completion or to satisfy
\(T_{\tau_S}(F^r)\subseteq F^{r+1}\), so that every iterated expression in
the completed tensor product converges.
For homogeneous maps use the declared convolution product and Hom
differential
\[
 (f\star g)(c)=
 \sum(-1)^{|g||c_{(1)}|}f(c_{(1)})g(c_{(2)}),
 \qquad
 d_{\mathrm{Hom}}f=d_Ef-(-1)^{|f|}f d_P.
\]
Prove the Maurer--Cartan equation
\[
 d_{\mathrm{Hom}}\tau_S+\tau_S\star\tau_S=0,
\]
write the induced twisted differential on
\(P_S\widehat\otimes B_S\) with all Koszul signs, and verify its square
directly.

\item If \(\tau_S\) is to come from tame symbols, construct its
chain-level factorization through a specified residue-unit group algebra
and an action of that algebra on \(B_S\).  A multiplicative reciprocity
identity alone is not a Maurer--Cartan proof.

\item For every inclusion \(S\subset T\), construct a chain map
\[
 j_{S,T}:(P_S\widehat\otimes B_S,D_{\tau_S})
 \longrightarrow(P_T\widehat\otimes B_T,D_{\tau_T})
\]
with strict or specified homotopy-coherent composition.  Prove that it
intertwines the twisting operators.  Only then form the filtered colimit
and its cohomology.  Any claimed concentration in degrees \(0,1,2\) must
be proved; three untwisted polynomial prime factors already have the
nonzero class \(\eta_p\wedge\eta_q\wedge\eta_r\) in degree three.

\item Relate a finite-place algebra to the Toeplitz carrier by an actual
representation on the dense space \(\bigoplus_{r\geq0}\mathbb C e_{p,r}\)
and an Euler derivation \(E_p\) satisfying
\[
 [N_p,\pi_p(a)]=\pi_p(E_pa)
\]
on a common invariant domain.  For the divisor algebra this part is
already supplied by \(\pi_p(u_p)=S_p\) and
\(E_p(u_p^r)=r\log p\,u_p^r\): direct evaluation on the finite basis span
gives \([N_p,S_p^r]=r\log p\,S_p^r\).
The additional arithmetic coefficient algebra must carry its own such
representation and comparison with this divisor model.
\end{enumerate}

No system with the stated arithmetic, archimedean-chiral, tame-symbol, and
transition properties is constructed here.  Until one is constructed, its
colimit and cohomology are specifications rather than arithmetic-chiral
invariants.
\end{openproblem}
